\section{Методология сквозной верификации на базе Verilator}

Классический цикл верификации цифровых узлов в САПР вендоров ПЛИС опирается на событийно-ориентированные симуляторы. В таких средах изменение логического уровня в любом узле порождает событие в глобальной очереди планировщика. При обработке непрерывных выборок радиосигналов на промежуточной частоте количество таких событий на каждый такт тактовой частоты исчисляется десятками тысяч, что приводит к критической деградации производительности моделирования.

Альтернативой выступает применение открытого инструмента Verilator \cite{verilator, snyder_verilator}, выполняющего статическую компиляцию синтезируемого подмножества SystemVerilog в оптимизированный C++ код. Verilator лишен классической очереди событий: проект преобразуется в детерминированный ориентированный граф логических операций, который транслируется в монолитный класс \texttt{Vtop}. Внутренние регистры и цепи проекта становятся штатными переменными в оперативной памяти хост-системы, а модельное время продвигается явным переключением тактового сигнала и вызовом метода \texttt{top->eval()}.

Прямая трансляция аппаратного описания в код на C++ принципиально меняет архитектуру совместного моделирования (HW/SW Co-Design). Обмен данными между синтезируемым ядром и верификационной обвязкой сводится к прямому чтению и записи полей C++ структуры, исключая накладные расходы на межпроцессное взаимодействие (IPC) и трансляцию типов через интерфейсы DPI-C или VPI. Это позволяет бесшовно подключать к аппаратному ядру произвольные внешние математические библиотеки, реализовывать ресурсоемкую статистическую обработку в нативной арифметике с плавающей точкой (\texttt{float}, \texttt{double}) и собирать весь верификационный комплекс в единый высокоскоростной бинарный исполняемый файл средствами стандартных компиляторов GCC или Clang.
